% !TeX root = ../papers/nc-p3.tex
\section{四つの基点を通る判別三次式と三平面の楕円曲線}
\label{sec:marked-cubic-linear-system}
\begin{theorem}\label{thm:marked-cubic-linear-system}
$G'\subset S=\mathrm{Sym}_0(V)$ を固定した有限混合印の対称三平面とし、
$\bigwedge^2G'\to\mathfrak{so}(V)$ は可逆とする。$G'^\perp$ の基底 $N_1,N_2$ が
定める二次式 $q_i(v)=v^tN_i v$ の完全交叉を $B$ とする。
$G'$ を固定した全ての相手三平面の混合行列式は $H^0(\mathcal I_B(3))$ に入る。
逆にこの六空間の\emph{任意の}非零三次式 $f$ を判別式に指定して、元の混合点を通る
幾何学的一般点で全次数正則・四辺有限・反対階数10の族を取れる。
$f=l_1q_1+l_2q_2$ の一意な表示で $l_1,l_2$ が独立なら、その一般点で $g=5$ である。

特に分裂点 $G_i=D_i$ では $B=\{(1,\pm1,\pm1)\}$ であり、
この四点を通る各 $f$ に対し、同型類の像が六次元、一般最小周期4、一般射影自己同型群が
自明な族となる。従って滑らかな三次曲線の全ての $j$ 不変量を含み、
台が一本の直線より大きいどの平面三次曲線も射影同値までこの構成で実現する。
\end{theorem}
\begin{proof}
二次式に共通因子があれば両方が同じ一次式 $v\cdot n$ を持ち、
跡零性から他の一次因子二つは $n^\perp$ を張る。従って $G'$ は全元が
$kn$ を保つ三平面となり、有限印の既約性に反する。二式は正則列で、飽和イデアルの
三次部分は $q_1,q_2$ の一次式倍の直和である。
$\omega_v(X,Y)=\operatorname{tr}([v]_\times[X,Y])$ の根基を生成する
$T_v=vv^t-(v^tv)I/3$ は $B$ 上スキームとして $G'$ に入る。各点でその一成分は単元なので、
任意の $F$ に対し $\det(\omega_v(G_a,F_c))$ は $B$ 上で零である。

この行列式写像の $F=-G$ での微分を計算する。$F$ の $G'$ 内での変動は微分で零になる。
$Z\in\mathrm{Hom}(k^3,S/G')$ が階数1で $Z_a=\eta_a Y$ なら、
$\ker\eta$ の二基底の $G$ による像を $H_1,H_2$ とする。
命題\ref{prop:marked-mixed-discriminant}のPfaffian展開により、一次係数は
$\omega_v(H_1,H_2)\, q_N(v)$ に比例し非零である。ここで $N$ は $G'$ と $Y$ を消す法線である。
実際行列式において一列だけを $Y$ に替える基底を使えるので係数は一次が全てである。
三交換子の独立性から $[\eta]\mapsto[\omega_v(H_1,H_2)]$ は射影線型同型であり、
$[Y]\mapsto[N]$ は二次元商での消去によりやはり射影線型同型である。
従って微分の像は全ての一次式と任意の鉛筆二次式の積を含む。その線型包は六次元で、
微分は六空間の同型である。特にこの同型は階数1の射影Segreを、
$l_1,l_2$ の従属するSegreに移す。

跡条件 $\sum[G_i,F_i]=0$ は $G'$ 内の係数で任意に補える。
よって固定 $G$ の十二次元の $F$ 座標からこの六空間への写像は混合点で滑らかであり、
指定した直線 $kf$ の逆像は同点で滑らかな七次元で、実際に直線に支配的である。
交代解のエタール復元と内部次数別の相対bar完全性を前節までと同様に適用する。
$l_1,l_2$ が独立なら階数1の変動ではこの直線方向に出られず、階数2の方向があるから
一般点で $g=5$ となる。

分裂点では対角跡零行列が法線鉛筆で、表示した四点を得る。$G'$ を固定して
順序付き基底も動かすと次元は $7+9=16$、$G'$ を保つ射影直交群は有限であり、
残余群の連結成分は次元10である。この群は判別式の射影値を保ち、印は有限なので
像次元は6となる。定理\ref{thm:marked-semisimple-period}の混合点内の基底変動はこの族でも取れ、
$kf$ の非零方向は必ず対角商で非零の変動を持つ。同定理の六列が独立な持上げ障害は
そのまま適用できる。最後に三次曲線の台から射影枠四点を選んで上の四点に移せばよい。
二直線に支持される場合にも各直線から二点を選べる。三重直線だけはこの枠を持たないが、
その射影型は等方一次式とその平方を定理\ref{thm:marked-prescribed-discriminant-as}で指定して実現できる。
この射影型の議論では曲線の座標を任意に取り替えており、指定計量との相対位置は固定していない。
\end{proof}

\begin{theorem}\label{thm:marked-elliptic-complement-planes}
分裂三平面 $G'=\langle D_1,D_2,D_3\rangle$ を固定し、滑らかな $E=V(f)$ が
上の四点を通るとする。$F'\in\mathrm{Gr}(3,S)$ について、混合行列式が $kf$ に入る
全三平面のスキームは、次数5の楕円正規曲線で $\mathrm{Pic}^2(E)$ と同型である。
そのうち行列式が恒等零になるのはちょうど五つの被約点で、
$G'$ と、四点 $v\in B$ について $kv$ を保つ跡零対称三平面である。
他は全て $g=5$ である。これは対称側の相手の\emph{部分空間}を尽くす。
順序付き基底の跡条件、三交換子の持上げとその特殊点のAS条件は別に課す。
\end{theorem}
\begin{proof}
$G'$ と $S$ の基底で $A(v)=(\omega_v(G_a,S_j))$ を三行五列の一次行列とする。
その三次小行列式は前定理により $H^0(\mathcal I_B(3))$ を張る。相手の三平面の
Pl\"ucker埋込み $\mathrm{Gr}(3,5)\subset\mathbf P^9$ において混合行列式は六成分の線型写像である。
中心の $\mathbf P^3$ とGrassmann多様体は表示した五点で横断する。
\path{computations/verify_mixed_cubic_pencil.py}が三つの非対角対称行列と二つの対角跡零行列を
基底に用い、各三平面のアフィンチャートの六座標に関する小行列式微分が
有理数上で階数6であることを検算する。Grassmann多様体の次数は5
（$2\times3$ 長方形の標準盤の数）である。線型切断を順に取ると成分の次数の和は
増えないため、横断的な孤立五点で全切断を尽くす。

$E$ 上の核束 $\mathcal U=\ker(S\otimes\mathcal O_E\xrightarrow A (G')^*\otimes\mathcal O_E(1))$
を取る。$B$ でも像は曲線上で局所自由なので核は部分束である。
$A$ の階数は $T_v\notin G'$ で3、$T_v\in G'$ で2である。
各基点で小行列式は $E$ 上の極大イデアルを生成する。実際六次元の三次式系が
二つの二次生成元の全一次倍を含み、その一方以上は滑らかな $E$ 上で一次に消える。
従ってこの写像の余核の長さは4で、$\deg\mathcal U=-5$、
$\det\mathcal U^*=\mathcal O_E(3)(-B)$ である。
$H^0(\mathcal U)$ の元は線型行列 $A(v)$ により曲線上恒等的に消える定数ベクトルで、
曲線は平面を張るから $G'$ の全元と交換しなければならない。既約性と跡零性からこれは零で、
種数1のRiemann--Rochにより $h^0(\mathcal U^*)=5$ である。

自然写像 $S^*\to H^0(\mathcal U^*)$ は同型である。実際その核に $a\ne0$ があれば、
$a\wedge S^*$ の射影三空間は可分解なPl\"ucker点のみからなり、全て $E$ 上の
$\bigwedge^2\mathcal U^*$ への制限が零である。この制限が零という条件は、
三平面 $F'$ の消去二平面を取り、上の線型写像で行列式が $kf$ に入る条件と一致する。
一般の $v\in E$ で二平面 $\mathcal U_v$ が $F'$ と交わる条件そのものであり、
小行列式は $\det\mathcal U^*$ の切断として基点因子を取り除いたものである。
従って先の射影三空間が $kf$ の逆像の $\mathbf P^4$ に入り、その超平面である中心とは
少なくとも平面で交わり、中心の有限性に反する。

$\mathcal U^*$ は安定である。次数3以上の線部分束があればその切断三つ以上の
二つずつの外積は全て $E$ 上で零となる。三次元切断部分の外積の射影平面は可分解点だけからなり、
上の $\mathbf P^4$ の中で中心と少なくとも直線で交わってしまう。
そこで任意の $M\in\mathrm{Pic}^2(E)$ に対し $\mathcal U^*\otimes M^{-1}$ は次数1の安定階数二束で、
双対に切断がなくRiemann--Rochから切断は一次元である。
得られる $M\hookrightarrow\mathcal U^*$ に零点があれば飽和の次数は3以上になるので部分束である。
$H^0(M)$ は基点自由な二平面で、これを $S^*$ 内に取れば対応する $F'$ が得られる。
逆に $kf$ の逆像のGrassmann点は二つの独立切断を与え、その評価は $E$ 上常に従属する。
像の飽和は独立切断を二つ持つ線束なので次数2以上で、安定性からちょうど2である。
従って元の $M$ が一意に戻る。

この対応はPicard曲線からの射である。普遍線束に対する上の切断空間は次元一定なので
局所的に束となり、二平面はGrassmann多様体へ正則に写る。逆像の台は以上より既約な曲線であり、
その正規化はこの全単射と標数零から $\mathrm{Pic}^2(E)$ である。
五超平面による所期次元の完全交叉はCohen--Macaulayで、$F'=G'$ で滑らかなので被約である。
$K_{\mathrm{Gr}(3,5)}=\mathcal O(-5)$ より随伴公式で算術種数1となり、
正規化の種数と一致するため滑らかである。次数は5であり、五基点のPl\"uckerベクトルは前掲の検算で階数4なので中心を張る。
曲線は中心外にも点を持つので全体の線型包は四空間である。
従って楕円正規五次曲線である。$g=4$ なら判別式は可約なので、
五基点を除いて $g=5$ に限る。
\end{proof}
